
Whether a base model's output probability distribution encodes information about which items will change their predicted answer after alignment is an open empirical question. If such a signal exists, it would enable item-level pre-deployment risk assessment: practitioners could identify evaluation questions most likely to receive different answers after alignment before any fine-tuning is applied.

The central object of study is the \textit{argmax flip} --- the event that a model's top-predicted answer for a given MCQ item changes after an alignment procedure. Among 14,042 MMLU test items, 12.5\% receive a different top prediction after DPO alignment of tulu-2-7b, with no change to the question content. The empirical question is whether properties of the base model's output distribution --- observable before alignment --- predict which items will flip.

The \textit{pre-alignment confidence margin} is defined as the z-scored difference between the base model's top-1 and top-2 log-probabilities over the MCQ answer options. This quantity captures the geometric distance from the item's implied answer distribution to the decision boundary at which the top-1 and top-2 answers tie. Items with small margins lie near this boundary; items with large margins lie far from it.

The hypothesis tested in this paper is that low-margin items are disproportionately exposed to alignment-induced boundary shifts, because alignment perturbations of bounded magnitude are sufficient to cross the boundary when the initial margin is small. This predicts a negative relationship between pre-alignment confidence margin and post-alignment argmax flip probability.

Three research questions are examined:

\textbf{RQ1.} Does pre-alignment confidence margin predict post-alignment argmax flip probability at above-threshold discriminative performance, and does this generalize across benchmark domains?

\textbf{RQ2.} Are alignment-induced logit perturbations geometrically structured --- specifically, non-isotropic --- as opposed to approximating random noise?

\textbf{RQ3.} Do DPO and SFT produce distinguishably different perturbation variance profiles across confidence quintiles?

This paper reports the following: (1) for a DPO-aligned tulu-2-7b model pair, the margin predictor achieves AUROC of 0.867--0.909 across three benchmarks with no domain adaptation; (2) logit-delta perturbations are non-isotropic under both DPO and SFT with anisotropy ratios of 2.90 and 4.58 respectively; (3) contrary to the directional prediction, DPO does not concentrate perturbation variance in low-confidence regions --- instead DPO variance increases monotonically with confidence (a null result for the directional hypothesis, with a novel finding about DPO's quintile profile). Limitations include the absence of PPO-aligned model pairs, primary evidence from a single model pair for the margin-flip result, and MCQ-specific scope.

